\documentclass{beamer}

\usetheme{Berlin}
\usecolortheme{Orchid}
\useoutertheme{miniframes}
\setbeamertemplate{navigation symbols}{}

\usepackage[utf8]{inputenc}
\usepackage[T1]{fontenc}
\usepackage{amsmath}
\usepackage{amssymb}
\usepackage{booktabs}
\usepackage{graphicx}
\graphicspath{{../images/}}

\usepackage{xcolor}
\usepackage{listings}
\usepackage{hyperref}
\usepackage{tikz}
\usetikzlibrary{arrows.meta,positioning,shapes.geometric,calc}

\lstset{
  language=Java,
  basicstyle=\ttfamily\small,
  keywordstyle=\color{blue},
  commentstyle=\color{gray},
  stringstyle=\color{teal},
  showstringspaces=false,
  columns=fullflexible,
  breaklines=true
}

% Per-slide source line (references shown on the slide itself).
\newcommand{\slidesources}[1]{%
  \vfill
  {\tiny\color{gray}\raggedright\sloppy\parbox{\linewidth}{Sources: #1}\par}%
}

% Unit-wise source strings for this subject (used by slides/notes).
% Path is relative to the lecture `latex/` folder (the compilation CWD).
% OOPS with Java (CSEG1044) — unit-wise sources
%
% These are short, student-facing reference strings intended to be used with:
%   - \slidesources{...}  (in beamer slides)
%   - \notesources{...}   (in lecture notes)
%
% Style inspiration: other subjects in My-Subjects/ use semicolon-separated sources
% and include an "accessed" date for web documentation.

\newcommand{\OOPSUnitISources}{Schildt, \textit{Java: A Beginner's Guide} (10e), McGraw-Hill (Java fundamentals + OOP basics).; Oracle Java Tutorials: Getting Started + Language Basics + Classes and Objects (accessed \today).; Java SE API docs (e.g., \texttt{String}, \texttt{Scanner}) (accessed \today).}

\newcommand{\OOPSUnitIISources}{Schildt, \textit{Java: The Complete Reference} (12e), McGraw-Hill (classes, inheritance, packages, interfaces).; Bloch, \textit{Effective Java} (3e), Addison-Wesley, 2018 (design choices; interfaces vs abstract classes).; Oracle Java Tutorials: Classes and Objects + Interfaces + Packages (accessed \today).; Java SE API docs (e.g., \texttt{Comparable}, \texttt{Runnable}) (accessed \today).}

\newcommand{\OOPSUnitIIISources}{Schildt, \textit{Java: The Complete Reference} (12e) (nested/inner classes, exceptions, threads, I/O).; Oracle Java Tutorials: Nested Classes + Exceptions + Concurrency + Basic I/O (accessed \today).; Java SE API docs (e.g., \texttt{Thread}, \texttt{AutoCloseable}) (accessed \today).}

\newcommand{\OOPSUnitIVSources}{Schildt, \textit{Java: The Complete Reference} (12e) (generics, lambdas, Swing, JDBC).; Bloch, \textit{Effective Java} (3e) (generics + lambdas).; Oracle Java Tutorials: Generics + Lambda Expressions + Swing + JDBC Basics (accessed \today).; Java SE API docs (e.g., \texttt{java.util.function}, \texttt{javax.swing}, \texttt{java.sql}) (accessed \today).}

\newcommand{\OOPSUnitVSources}{Schildt, \textit{Java: The Complete Reference} (12e) (Collections Framework + wrapper classes).; Oracle Java docs: Collections Framework overview (accessed \today).; Java SE API docs (\texttt{java.util} collections; wrapper classes in \texttt{java.lang}) (accessed \today).}

\newcommand{\OOPSUnitVISources}{Git SCM: \textit{Pro Git} (2e) (version control workflow), accessed \today.; GitHub Docs (repositories, issues, pull requests), accessed \today.; Oracle Java Tutorials: Swing + JDBC (accessed \today).; JUnit 5 User Guide (accessed \today).; Apache Maven Project: Getting Started (accessed \today).; Gradle User Manual: Getting Started (accessed \today).; UML class diagrams: UML 2.x overview/spec (accessed \today).}

\newcommand{\OOPSMiscWhyInterfacesSources}{Schildt, \textit{Java: The Complete Reference} (12e) (interfaces).; Bloch, \textit{Effective Java} (3e) (prefer interfaces where appropriate).; Oracle Java Tutorials: Interfaces (accessed \today).; Java SE API docs (e.g., \texttt{Comparable}, \texttt{Runnable}) (accessed \today).}



\title[OOPS with Java]{Object Oriented Programming (CSEG1044)}
\subtitle{Unit 01 -- Lecture 04: Control Flow + Arrays + Strings}
\author{Tofik Ali}
\institute{School of Computer Science, UPES Dehradun}
\date{\today}

\begin{document}

\begin{frame}[fragile]
  \titlepage
  \vspace{0.5em}
  \slidesources{\OOPSUnitISources}
\end{frame}

\begin{frame}{Quick Links}
  \centering
  \hyperlink{sec:overview}{\beamerbutton{Overview}}\hspace{0.6em}
  \hyperlink{sec:concepts}{\beamerbutton{Key Topics}}\hspace{0.6em}
  \hyperlink{sec:demo}{\beamerbutton{Demo}}\hspace{0.6em}
  \hyperlink{sec:summary}{\beamerbutton{Summary}}
  \slidesources{\OOPSUnitISources}
\end{frame}

\begin{frame}{Agenda}
  \tableofcontents
  \slidesources{\OOPSUnitISources}
\end{frame}

\section{Overview}
\label{sec:overview}

\begin{frame}{Learning Outcomes}
  \begin{itemize}[<+->]
    \item Use \texttt{if/else}, \texttt{switch}, and loops to control program flow.
    \item Create/traverse 1D and 2D arrays and solve simple array problems.
    \item Explain \texttt{String} immutability and use \texttt{StringBuilder} when needed.
  \end{itemize}
  \slidesources{\OOPSUnitISources}
\end{frame}

\section{Key Topics}
\label{sec:concepts}

\begin{frame}{Unit 01: Introduction to OOPs}
  \begin{itemize}[<+->]
    \item Lecture focus: Control Flow + Arrays + Strings
  \end{itemize}
  \slidesources{\OOPSUnitISources}
\end{frame}

\begin{frame}{Today: Roadmap}
  \begin{itemize}[<+->]
    \item Program control statements (selection + loops)
    \item Arrays (1D/2D), types of arrays, common patterns
    \item Strings (immutability, common methods), `StringBuilder` (need-based)
  \end{itemize}
  \slidesources{\OOPSUnitISources}
\end{frame}

\begin{frame}{Control flow (quick)}
  \begin{itemize}[<+->]
    \item Selection: \texttt{if/else}, \texttt{switch}
    \item Loops: \texttt{for}, \texttt{while}, \texttt{do-while}
    \item Useful keywords: \texttt{break}, \texttt{continue}
  \end{itemize}
  \begin{block}{Common pitfalls}
    \begin{itemize}
      \item Off-by-one errors in loops
      \item Missing \texttt{break} in \texttt{switch} (unintended fall-through)
    \end{itemize}
  \end{block}
  \slidesources{\OOPSUnitISources}
\end{frame}

\begin{frame}{Arrays (1D/2D)}
  \begin{itemize}[<+->]
    \item 1D: \texttt{int[] a = new int[n];} and \texttt{a.length}
    \item Traversal: \texttt{for (int i=0; i<a.length; i++)} or \texttt{for (int x: a)}
    \item 2D: \texttt{int[][] m = new int[r][c];}
    \item Jagged arrays: each row can have different length (use \texttt{m[i].length})
  \end{itemize}
  \slidesources{\OOPSUnitISources}
\end{frame}

\begin{frame}{Strings (must-know)}
  \begin{itemize}[<+->]
    \item \texttt{String} is \textbf{immutable}: methods return new strings.
    \item Compare content with \texttt{equals()}, not \texttt{==}.
    \item Repeated concatenation in loops is slow $\rightarrow$ use \texttt{StringBuilder}.
  \end{itemize}
  \slidesources{\OOPSUnitISources}
\end{frame}


\section{Demo / Example}
\label{sec:demo}

\begin{frame}[fragile]{Example: grades (arrays + control flow) (1/2)}
\begin{lstlisting}
import java.util.Scanner;

public class GradeStats {
    public static void main(String[] args) {
        Scanner sc = new Scanner(System.in);
        System.out.print("Enter number of students: ");
        int n = sc.nextInt();

        int[] marks = new int[n];
        for (int i = 0; i < n; i++) {
            System.out.print("Marks for student " + (i + 1) + ": ");
            marks[i] = sc.nextInt();
        }
\end{lstlisting}
  \slidesources{\OOPSUnitISources}
\end{frame}

\begin{frame}[fragile]{Example: grades (arrays + control flow) (2/2)}
\begin{lstlisting}
// continued...
        int max = marks[0], sum = 0;
        for (int x : marks) { sum += x; if (x > max) max = x; }
        double avg = sum / (double) n;
        System.out.printf("Max=%d, Avg=%.2f%n", max, avg);
        sc.close();
    }
}
\end{lstlisting}
  \slidesources{\OOPSUnitISources}
\end{frame}

\begin{frame}{Mini Exercises (2--3 mins)}
  \begin{enumerate}[<+->]
    \item Array of length 10: what is the last valid index?
    \item Which compares string content: \texttt{==} or \texttt{equals}?
    \item Which loop runs at least once: \texttt{while} or \texttt{do-while}?
  \end{enumerate}
  \slidesources{\OOPSUnitISources}
\end{frame}


\section{Summary}
\label{sec:summary}

\begin{frame}{Key Takeaways}
  \begin{itemize}[<+->]
    \item Use selection + loops to express logic clearly (avoid off-by-one bugs).
    \item Arrays store many values of the same type; \texttt{length} guides safe iteration.
    \item Strings are immutable; use \texttt{equals()} and \texttt{StringBuilder} when needed.
  \end{itemize}
  \slidesources{\OOPSUnitISources}
\end{frame}

\begin{frame}{Checkpoint Questions}
  \begin{enumerate}[<+->]
    \item How do you avoid off-by-one errors in array loops?
    \item Why should we use \texttt{equals()} for strings?
    \item When is \texttt{StringBuilder} preferred over \texttt{String} concatenation?
  \end{enumerate}
  \slidesources{\OOPSUnitISources}
\end{frame}

\begin{frame}{Next Unit Preview}
  \textbf{Next:} Unit 02 -- Classes, Objects, and Encapsulation
  \vspace{0.6em}
  \begin{itemize}[<+->]
    \item We will start writing proper classes with private fields and methods.
    \item Bring one example of a real-world entity you want to model as a class.
  \end{itemize}
  \slidesources{\OOPSUnitISources}
\end{frame}

\end{document}
